
\Slide{1}{函数与对应的入口}
本讲把“唯一输出”的函数推广为“输出集合”的对应。先学集合元素与子集的区别：$y\in\Phi(x)$，而不是把集合 $\Phi(x)$ 当成数。

\Slide{2}{本讲的位置}
集合运算用于处理像与逆像，二元运算用于建立实数代数，后面的连续性则通过逆像描述。先把符号翻译为中文句子，再证明集合等式。

\Slide{3}{函数是什么}
$f:S\to T$ 要求每个 $x\in S$ 恰有一个输出。定义域是 $S$，陪域是 $T$。例如 $f(x)=x^2$ 从实数到实数是函数，虽然不同输入可能有相同输出。

\Slide{4}{函数的图}
图是所有 $(x,f(x))$ 的集合。证明某集合是函数图，要检查每个输入存在且仅有一个对应点；“画得像曲线”不构成证明。

\Slide{5}{像与满射}
$y\in f(A)$ 等价于存在 $x\in A$ 使 $f(x)=y$。满射是 $f(S)=T$；改变陪域可改变满射性质，例如平方函数到 $\R$ 不满，到 $[0,\infty)$ 满。

\Slide{6}{像保持并集}
$y\in f(\bigcup A_i)\iff\exists i\exists x\in A_i:f(x)=y\iff y\in\bigcup f(A_i)$。交集只得包含：同一个 $x$ 属于所有 $A_i$ 才可直接取像；反向可能由不同原像产生同一个 $y$。

\Slide{7}{指标写法与反例}
本页把上页集合族改为指标族，证明不变。常值函数 $f(1)=f(2)=0$ 给出 $f(\{1\}\cap\{2\})=\varnothing$ 而像的交为 $\{0\}$。反例必须把左右两边都算出。

\Slide{8}{逆像不用反函数}
$f^{-1}(B)=\{x:f(x)\in B\}$ 对任何函数都定义。例如平方函数下 $\{1\}$ 的逆像是 $\{-1,1\}$。逆像是集合，不要求 $f$ 可逆。

\Slide{9}{逆像与逻辑}
$x\in f^{-1}(\bigcap B_i)\iff\forall i,f(x)\in B_i\iff x\in\bigcap f^{-1}(B_i)$。把全称改为存在即证并集公式。差集公式来自“属于 $B$ 且不属于 $B'$”。

\Slide{10}{逆像公式的另一写法}
指标族与集合族表达同一事实，重写上页的三个等价链即可。补集公式是 $f^{-1}(T\setminus B)=S\setminus f^{-1}(B)$，注意左右补集的全集不同。

\Slide{11}{单射与双射的区别}
纠错：每个 $y\in T$ 恰有一个原像表达双射，不只是单射。单射仅要求 $f(x)=f(x')\Rightarrow x=x'$；陪域中的某些点可以没有原像。双射同时满足单射和满射。

\Slide{12}{延拓与限制}
限制缩小定义域但保持原来输出，延拓扩大定义域并在旧定义域与原函数一致。例：实数平方函数限制到非负实数后成为双射。检查延拓时应逐个核对旧输入，而不是只比较公式名称。

\Slide{13}{坐标投影}
$\pi_i(x_1,\ldots,x_N)=x_i$。投影会丢失其他坐标，一般不单射。以后证明向量收敛可逐坐标进行，因为坐标数是有限的。

\Slide{14}{对应与经济选择}
$\Phi(x)\subseteq T$ 可以含多个元素，例如预算下所有可行消费。课件有时默认值非空，有时允许空值；每道题应说明采用哪种约定。

\Slide{15}{有效定义域}
$S^*=\{x:\Phi(x)\ne\varnothing\}$，值域为所有输出的并。单值且处处非空时可以视为函数。不要把“对应的值域”等同于一开始指定的陪域。

\Slide{16}{集合的对应像}
$\Phi(A)=\bigcup_{x\in A}\Phi(x)$，即从 $A$ 中至少一个输入可以达到的输出。若 $A$ 为空，像按约定为空。集合族的像则是由这些像组成的新集合族，层次不同。

\Slide{17}{对应像的运算}
单调性：$A\subseteq B$ 时任一 $y\in\Phi(A)$ 有见证 $x\in A\subseteq B$，故属于 $\Phi(B)$。并集交换存在量词得等式；交集只有包含。补集结论必须相对实际值域，否则未达到的输出没有原像可选。

\Slide{18}{补集结论的证明}
在 $T=\Phi(S)$ 下，若 $y\notin\Phi(A)$，从满值域选 $x\in S$ 使 $y\in\Phi(x)$；该 $x$ 不在 $A$，故 $y\in\Phi(S\setminus A)$。反向需不同输入的输出集不相交，否则同一输出可同时由 $A$ 内外产生。

\Slide{19}{对应的图}
$G(\Phi)=\{(x,y):y\in\Phi(x)\}$。竖直截面等于值集，可以有多个点。104中二元关系也能看成这类有序对集合。

\Slide{20}{二元运算}
运算是 $S\times S\to S$ 的函数，必须封闭且单值。正整数的减法不封闭，因此不是正整数上的二元运算；加法和乘法是。

\Slide{21}{交换律与结合律}
交换律交换两个数位置，结合律改变三个数的括号。幂运算不交换：$2^3\ne3^2$；不结合：$(2^2)^3=64$ 而 $2^{(2^3)}=256$。一个反例足以否定全称命题。

\Slide{22}{分配律}
乘法对加法分配：$a(b+c)=ab+ac$。加法不对乘法分配，例如 $1+(2\cdot3)=7$，但 $(1+2)(1+3)=12$。运算顺序需逐层读取。

\Slide{23}{对应逐点运算}
并、交、积均先固定输入 $x$ 再对值集运算。交值可能为空，即使两个原值都非空，因此不能自动保留非空值条件。

\Slide{24}{对应的复合}
$(\Psi\circ\Phi)(x)=\bigcup_{y\in\Phi(x)}\Psi(y)$，表示经某个中间点能到达的输出。逐点交的像一般小于像的交，因为后者可使用不同输入。积的像也可能只含同步产生的坐标对。

\Slide{25}{常对应与恒等对应}
若 $\Delta(x)=C$，对非空 $A$ 有 $\Delta(A)=C$，所以 $\bigcup_{x\in A}(C\cap\Phi(x))=C\cap\Phi(A)$；空 $A$ 时两边均为空。恒等对应的输出是单点集 $\{x\}$。

\Slide{26}{复合结合与可达性}
复合结合来自交换“存在中间点”的分组。传递闭包 $\widehat\Phi$ 收集走零步、一步、两步等可达点。再走有限步仍为有限步，故 $\widehat\Phi(\widehat\Phi(A))=\widehat\Phi(A)$；零步给扩张性，见证包含给单调性。

\Slide{27}{半单值及单射}
半单值指每个值集至多一个元素。交与积保留该性质，因为从每个因子至多各取一个元素。若某个因子在不同输入的值集不相交，则积中同一点不可能同时来自两个输入；交同理。

\Slide{28}{逆对应}
$\Phi^{-}(y)=\{x:y\in\Phi(x)\}$。若原对应每点恰一输出，则两个不同 $y$ 的逆值不相交。若原对应不同输入值集不相交，则每个 $y$ 至多有一个原像。纠错：证明后半必须使用“不同输入的值集不相交”，不能从单射推出同一输入有唯一输出。

\Slide{29}{下逆像}
$x\in\Phi^{-}(B)\iff\exists y\in B:y\in\Phi(x)$。交换有序对两坐标再交换回来不变，故逆对应的逆是原对应。集合逆像的“存在”含义对后面下半连续极重要。

\Slide{30}{上逆像}
$\Phi^+(B)=\{x\in S^*:\Phi(x)\subseteq B\}$ 表示所有允许输出都在 $B$。非空值时“全部在内”推出“至少一个在内”，所以 $\Phi^+(B)\subseteq\Phi^-(B)$。

\Slide{31}{补集对偶}
在有效定义域 $S^*$ 内，非“全部在 $B$”等价于“有一个在 $T\setminus B$”。因此 $S^*\setminus\Phi^+(B)=\Phi^-(T\setminus B)$。允许空值却取整个 $S$ 作补集会破坏课件公式。

\Slide{32}{并集的上下逆像}
下逆像保持并：有输出在 $B_1\cup B_2$ 当且仅有输出在至少一个集合。上逆像仅有包含，因为所有输出可能分散在两集合。例值集 $\{0,1\}$，分别取 $B_1=\{0\},B_2=\{1\}$ 即为严格包含。

\Slide{33}{数列是函数}
数列是 $\N\to S$ 的函数，$x_n$ 是第 $n$ 项，允许重复。集合不记录顺序及重复，数列记录；集合 $\{1,2\}$ 与交替数列 $1,2,1,2,\ldots$ 不能混为一谈。

\Slide{34}{可数与有限乘积}
按本课约定可数包括有限集。$\N^2$ 按坐标和依次列出：$(1,1),(1,2),(2,1),\ldots$，每个数对有限步后出现。归纳得有限多个可数集的积可数；这不意味着可数无穷积也可数。

\Slide{35}{Cantor 对角证明}
假设 $f:A\to\mathcal P(A)$ 满射。定义 $B=\{a\in A:a\notin f(a)\}$，由满射存在 $b$ 使 $f(b)=B$。于是 $b\in B\iff b\notin B$，两种可能都矛盾。关键是 $B$ 依赖于假设中的整张列表，必与每个列出集合在其对应位置不同。
